% !TEX root = ../thesis.tex
% Chapter 1: Introduction

\label{ch1:intro}

\section{Biomolecular structure as a problem of distributions}
\label{ch1:sec:biomolecular-distributions}

Historically, structural biology focused on the native structure---the dominant, stable fold of a biomolecule. At the dawn of the field, this description was apt. The early structures of DNA, myoglobin, and tRNA, solved from the 1950s through the 1970s, were resolved by X-ray diffraction as well-folded structures \citep{Watson1953DNA,Kendrew1958Myoglobin,Kim1974tRNA}. As the number of solved molecular structures grew, determining the native structure became the central goal of experimental structure determination and, later, of computational prediction.

Concurrently, theorists sought first-principles descriptions of biomolecular folding. For RNA, dynamic-programming approaches and nearest-neighbor energetic rules showed that secondary structure could be predicted from an explicit decomposition of interaction energies \citep{Nussinov1980FastAlgorithm,Mathews1999ThermodynamicParameters}. More generally, energy-landscape and folding-funnel pictures described folding as motion on a thermodynamic surface whose minima correspond to stable structures \citep{Bryngelson1987ProteinFolding,Wales2018EnergyLandscapes}. Spin-glass-inspired analyses further shed light on how heterogeneous interactions generate ordered, disordered, and metastable phases \citep{Bryngelson1987ProteinFolding,Pagnani2000Glassy}. The common insight is that many biomolecules have multiple accessible configurations, each with an associated thermodynamic weight, rather than a single, dominant fold.

From the 1980s onward, structural heterogeneity was increasingly confirmed experimentally: methods such as NMR spectroscopy, single-molecule FRET, and cryo-electron microscopy revealed structural variation that is not explained by a single static model of the structure \citep{Purcell1946NMR,Roy2008SingleMoleculeFRET,Kuhlbrandt2014ResolutionRevolution}. The goals of molecular modeling expanded beyond identifying one dominant structure, shifting toward characterization of structural ensembles.

The equilibrium probability distribution of statistical physics is the natural description: it assigns each microstate, or configuration, a Boltzmann weight. For a 3D molecular structure, a microstate is a set of atomic coordinates; for a spin model, it is a configuration of spin variables. In either case, the equilibrium distribution takes the form \citep{McQuarrie1976StatisticalMechanics}

\begin{equation}
  \label{thesis_intro:boltzmann}
  p_{\mathrm{eq}}(\mathbf{x};\beta)=\frac{e^{-\beta U(\mathbf{x})}}{Z(\beta)}\quad\mathrm{where}
  \quad
  Z(\beta)=\int e^{-\beta U(\mathbf{x})}\,d\mathbf{x}.
\end{equation}

The numerator in Eq. \ref{thesis_intro:boltzmann} is the Boltzmann weight. It describes how the probability of microstate $\mathbf{x}$ depends on its energy $U(\mathbf{x})$ for a system maintained by an environment at inverse temperature $\beta=1/k_BT$. At low temperatures, low-energy configurations dominate, while at higher temperatures the distribution flattens and high-energy states become accessible. The conceptual picture is intuitive, but estimating $p_\mathrm{eq}$ numerically is difficult in all but the most trivial cases. With an accessible energy function and temperature in hand, $p_\mathrm{eq}$ appears to be within reach. The obstacle is the normalization constant $Z$, called the partition function, whose evaluation requires an integral over a typically high-dimensional configuration space. Computing $Z(\beta)$ is of interest across physics, biology and chemistry, since the full thermodynamics of a system is encoded in the scalar-to-scalar map, a far simpler object than $p_{\mathrm{eq}}(\mathbf{x};\beta)$. However, practical, computationally tractable methods for obtaining $Z(\beta)$ proved to be elusive due to the so-called sampling problem.

Levinthal's paradox illustrates the severity of the problem: a 100-residue polypeptide has $10^{47}$ configurations, assuming each residue has 3 conformational states \citep{levinthal1969fold}. The astronomical configuration space makes enumeration or brute-force exploration impossible. Analytical methods from statistical mechanics for estimating $Z$, such as the replica trick, are useful, but they typically require assumptions that are violated by real biomolecular systems \citep{Parisi1997Replica}. Today's commonly used biomolecular energy functions like AMBER and CHARMM are necessarily complex because they must account for bonded interactions, electrostatics, van der Waals interactions, and solvent effects \citep{Wang2004GAFF,Vanommeslaeghe2010CGenFF}.

Another complication arises when one does not work with a fully atomistic description. In practice, biomolecular structure is often represented in terms of reduced variables or collective coordinates---for example, selected dihedral angles, radii of gyration, or contact distances. These reduced descriptions still follow equilibrium probability distributions, but the effective energy function that generates them is often not straightforwardly available. In such cases, the distribution must be learned from samples alone, without access to the energy function that generated them.

\section{Sampling equilibrium states of physical systems}
\label{ch1:sec:sampling-equilibrium-states}

The previous section identified the Boltzmann distribution and the partition function as central objects in statistical mechanics. However, most practical calculations do not begin by evaluating the partition function directly. Instead, they begin by generating samples from the equilibrium distribution. If representative configurations can be obtained, then structural properties, thermodynamic averages, and free-energy differences can be estimated from those samples. The central computational problem is therefore not only how to evaluate the partition function, but how to explore the relevant regions of configuration space in the first place.

One major class of sampling methods is Markov chain Monte Carlo (MCMC). MCMC constructs a sequence of configurations whose stationary distribution is the desired equilibrium distribution. In a typical MCMC algorithm, each step consists of a proposal and an acceptance decision. The proposal step generates a candidate configuration from the current one. The acceptance step then decides whether the chain moves to the proposed configuration or remains where it is. The acceptance rule is chosen so that, over many steps, the chain samples the target distribution \citep{metropolis1953equation,hastings1970monte}.

A central advantage of MCMC is that it does not require the normalization constant of the target distribution. Acceptance probabilities depend on ratios of probabilities, so the partition function is never required. This makes MCMC one of the most important tools in computational statistical physics: it allows equilibrium samples to be drawn from distributions with unknown or intractable normalization \citep{metropolis1953equation,hastings1970monte}.

At the same time, MCMC dynamics are not usually physical dynamics. The ``time'' in a Markov chain is algorithmic time: one counts MCMC steps, sweeps, or accepted moves, not seconds, minutes, or hours. A Markov chain may move between configurations in ways that are statistically valid but physically unrealistic. For this reason, MCMC is naturally suited to estimating equilibrium averages, but it does not generally provide kinetic information \citep{allen2017computer}.

The practical difficulty in MCMC is in designing the proposal step. Small local moves, such as flipping one spin or slightly perturbing one torsion angle, tend to be accepted often, but explore configuration space slowly. Large global moves can decorrelate samples more efficiently, but they are also more likely to propose high-energy configurations and be rejected. An effective MCMC implementation must simultaneously maintain high acceptance probabilities while exploring the configuration space efficiently.

A second major class of sampling methods is molecular dynamics (MD). Rather than proposing discrete transitions between configurations, MD simulates microscopic equations of motion. Given an energy function, forces are computed from the energy landscape and used to propagate the system forward in time \citep{verlet1967computer,allen2017computer}. In atomistic molecular simulations, the coordinates are atomic positions and the energy function is supplied by a molecular mechanics force field \citep{Wang2004GAFF,Vanommeslaeghe2010CGenFF}. With appropriate thermostats and barostats, MD simulations can be configured to sample particular thermodynamic ensembles, including the canonical and isothermal-isobaric ensembles \citep{nose--hoover,hoover,Bussi_Parrinello_thermostat,Parrinello_Rahman_dynamics}.

Like MCMC, performing MD simulations does not require the partition function directly. But unlike MCMC, MD simulations produce trajectories with an associated physical time variable. This makes them especially useful for the study of biomolecules, where the same simulation may provide both equilibrium structural information and insight into conformational dynamics. A trajectory of a protein, RNA, or ligand-receptor complex can, for instance, reveal which conformations are thermodynamically accessible and how the system transitions between them \citep{allen2017computer}.

The cost is that MD must resolve the fastest motions in the system. Bond vibrations and other high-frequency modes restrict the integration timestep to femtoseconds. At a timestep of one femtosecond, a microsecond of simulation requires one billion integration steps. A microsecond may be enough to observe a fast conformational rearrangement, but many biologically relevant events occur on longer timescales of seconds, hours, and days \citep{allen2017computer}.

MCMC and MD both face similar sampling limitations. Inefficient proposals in MCMC lead to rejection or slow mixing, while MD trajectories may remain trapped in local basins because barrier-crossing events are rare. In either case, the total number of sampled configurations can be misleading, since they are not statistically independent.

These limitations motivate enhanced sampling and free-energy methods. Some approaches modify the temperature or ensemble, some bias the dynamics along selected collective variables, and others reweight samples after they have been collected \citep{kirkpatrick1983optimization,SimulatedTempering,REMD,torrie1977nonphysical,metadynamics-review,abrams2014enhanced,wham,mbar,gallicchio2005temperature}. Despite their differences, they all address the same challenge of sampling high-dimensional distributions. The next section summarizes classical approaches developed to overcome this sampling problem and estimate partition functions or free-energy differences in realistic systems.


\section{Classical approaches to estimating partition functions and free energies}
\label{ch1:sec:classical-free-energy}

The difficulty introduced by Eq. \ref{thesis_intro:boltzmann} is not new. Much of computational statistical mechanics has been devoted to estimating partition functions, ratios of partition functions, or the free-energy differences derived from them. Since the absolute value of $Z(\beta)$ is rarely accessible for realistic molecular systems, most practical methods target relative quantities. For example, the free energy $F(\beta)=-\beta^{-1}\log Z(\beta)$ is discussed only in terms of free-energy difference between thermodynamic or conformational states. Classical approaches tend to focus either on improving the exploration of configuration space or extracting more accurate estimates of free-energy differences from samples that have already been collected.

Free-energy perturbation, Bennett acceptance ratio, thermodynamic integration, and targeted free-energy perturbation are the longest-standing route to computing free-energy differences \citep{zwanzig-fep,bar,TI,TFEP}. These methods apply when the relevant energy functions are known and when samples from one thermodynamic or conformational state overlap with the other. In practice, the overlap requirement is the central difficulty. Exponential averages can be dominated by rare samples, and free-energy estimates can converge slowly when the reference and target distributions differ substantially.

One family of methods improves sampling by changing the temperature or ensemble of the system. Simulated annealing raises and then lowers the temperature to help a system escape local minima \citep{kirkpatrick1983optimization}. Simulated tempering and replica exchange molecular dynamics extend this idea by allowing configurations to move among multiple temperatures or thermodynamic ensembles \citep{SimulatedTempering,REMD}. These {\it tempering} methods accelerate sampling by allowing high-temperature replicas to cross barriers, while maintaining low-temperature simulations to recover the free energy at a temperature of interest. However, the replicas must have similar energies for temperature exchanges to occur frequently enough to accelerate sampling. For large solvated biomolecules, maintaining this overlap can require many closely spaced replicas, which increases computational cost and limits scalability.

A second class of methods biases the dynamics along selected collective variables. Umbrella sampling, metadynamics, accelerated molecular dynamics, and related approaches modify the energy landscape so that rare transitions become frequent \citep{torrie1977nonphysical,metadynamics-review,abrams2014enhanced}. When the relevant slow coordinates are known, these methods can efficiently reconstruct free-energy profiles along those variables. The limitation is that the choice of collective variables is a strong modeling assumption. Many biomolecular systems do not evolve along one or two obviously meaningful coordinates. RNA conformational landscapes, in particular, can involve many competing base-pairing, stacking, tertiary-contact, and solvent-mediated interactions. In such cases, biasing a small set of variables may accelerate some motions while leaving other important barriers unresolved.

A final set of methods reconstructs thermodynamic information by reweighting samples after simulation. Weighted histogram analysis, multistate Bennett acceptance ratio, and related histogram-reweighting methods combine data from biased simulations or multiple thermodynamic states to estimate free-energy differences and, indirectly, ratios of partition functions \citep{wham,mbar,gallicchio2005temperature}. These approaches are especially valuable when samples have been collected at different temperatures or under different biasing potentials. Their accuracy, however, depends on sufficient overlap between the sampled distributions. If important regions of configuration space are missing, post-processing cannot recover their statistical weights.

The breadth of these approaches reflects the difficulty of the problem they aim to solve. Biomolecular free-energy landscapes are high-dimensional, rugged, and often only partially observed by experiment. Some methods require carefully chosen collective variables, while others require energy evaluations or overlap among sampled ensembles, and still others become expensive as system size increases. There is therefore no universal classical solution to estimating $Z(\beta)$, $\log Z(\beta)$, or the equilibrium distribution itself. This motivates the modern generative-model perspective: instead of selecting a small number of coordinates or relying only on reweighting, one can attempt to learn a map between a simple prior distribution and the complex molecular ensemble. The next section introduces this class of models and explains their utility for predicting molecular ensembles.

\section{Generative models for prediction of molecular ensembles}
\label{ch1:sec:generative-models}

The previous sections identify the equilibrium distribution, and the partition function that normalizes it, as central objects in biomolecular structure prediction. Classical methods improve sampling or estimate free-energy differences under specific assumptions, but they do not eliminate the need to model high-dimensional, multimodal distributions from sparse data. This problem has also attracted sustained interest in machine learning, where probabilistic generative models were developed to learn complicated probability distributions directly from data.

Probabilistic generative models (PGMs) begin with access to configurations drawn from an intractable target distribution. These may be data generated by experiment or simulation. A PGM learns a transformation between the samples from the target distribution and a simple prior distribution, conventionally a Gaussian. Generative models parameterize this transformation in many ways, but the end goal is always the same: to generate samples from the intractable target by sampling the simple prior and to evaluate the probability of data under the target distribution.

Normalizing flows were among the earliest successful realizations of this idea \citep{Dinh2017RealNVP}. They build the generative map as a composition of invertible transformations, so that the likelihood of a sample can be computed exactly through the change-of-variables formula. Exact evaluation of the likelihood is important because it permits training by likelihood maximization. However, normalizing flows are limited in expressivity, since the invertibility and likelihood calculation are made possible by restricting the mathematical operations the network performs. Moreover, it was discovered empirically that normalizing flows struggle to transform between a multimodal target distribution and a unimodal prior.

A major advance was to move the invertibility condition from the network architecture to an ordinary differential equation. In continuous normalizing flows, a neural network learns a velocity field that transports particles from Gaussian-distributed initial positions to the target distribution at a final time \citep{Grathwohl2018FFJORD}. Importantly, the likelihood remains accessible by integrating the divergence of that velocity field along the trajectory of the generative process. This continuous-time formulation is more flexible than discrete flows, but still struggles to generate multimodal distributions.

Diffusion models address this difficulty by employing a stochastic generative process. Inspired by non-equilibrium thermodynamics, they define a forward process that progressively destroys structure by adding noise to data, and a reverse process that reconstructs data by removing noise \citep{SohlDickstein2015NonEq,Ho2020DDPM,Song2021ScoreBased}. Because the forward process is stochastic, the learned reverse process can reconstruct disconnected metastable states without the mode-separation problems that limit deterministic models. The broader continuous-time literature now connects diffusion models, flow-matching variants, and related methods through a shared ODE--SDE--PDE perspective on probabilistic transport.

These developments have directly impacted structural biology. AlphaFold 3, for example, incorporates a diffusion model as part of a larger AI system. The success of AlphaFold 3 and contemporary models revealed another limitation. In many problems in molecular science, the main bottleneck is no longer the expressive power of the architecture, but the amount and quality of the available data. Generative models do not extrapolate reliably to entirely new classes or modes; they are only as good as the data on which they are trained. Protein structure prediction benefited from an unusually rich data regime. Other biomolecular structure-prediction problems, such as RNA structure prediction, have far less training data. Progress with limited high-quality molecular data therefore depends on inductive biases that promote generalization in AI and deep-learning models.

\section{RNA: A perfect storm of conformational heterogeneity, limited data, and environmental nuance}
\label{ch1:sec:rna-limited-data}

RNA structure prediction is a recurring topic throughout this thesis because it lies near the limit of what a purely data-driven generative model can be expected to learn. The difficulty is not only that fewer RNA structures are available than protein structures, but that the missing data are a consequence of the physics of RNA itself. RNA molecules are often more conformationally flexible than their protein counterparts: the ribose-phosphate backbone has six rotatable degrees of freedom per nucleotide. The large state space conferred by the flexible backbone is stabilized by base-pairing and -stacking interactions that generally give rise to multiple alternative states. In contrast, the protein backbone is frequently summarized by two main backbone dihedral angles per residue. This additional structural freedom makes RNA ensembles broader, more heterogeneous, and more difficult to summarize by a single representative conformation \citep{brion1997hierarchy,Wadley2007RNApseudotorsional}.

The same flexibility that makes RNA functionally rich also contributes to the scarcity of experimentally resolved structures. High-resolution structure-resolving methods such as X-ray crystallography, NMR spectroscopy, and cryo-electron microscopy are most straightforward to interpret when a molecule occupies a relatively coherent structural state. For flexible RNAs, the measured signal can instead be a superposition of conformations. Rather than resolving one clean structural model, the experiment may report an average, a partial model, or no high-resolution structure at all. The result is a limited number of solved structures that are biased toward RNAs that are experimentally tractable \citep{berman2000protein,bowman2024alphafold}.

RNA structure is also sensitive to environmental context. Ions such as sodium, potassium, and magnesium can reshape the folding landscape by neutralizing the negatively charged backbone or by stabilizing specific tertiary contacts. Some RNAs deliberately use this environmental sensitivity as the basis of function. Riboswitches regulate gene expression by adopting different structural states in response to metabolites or transition metals. Viral RNAs likewise encode structural elements whose stability and function depend on host context; dengue virus RNA elements, for example, have been studied at mosquito and mammalian temperatures. More generally, both coding and noncoding RNAs are frequently bound by protein partners or other nucleic acids, and the presence of those partners can change the relevant molecular structure \citep{yamagami2021functional,huang2025some,Villordo2015DengueRNA}.

For these reasons, RNA is not merely another version of protein structure prediction. It is a problem for which limited data, ensemble descriptions, and environmental factors compound one another. A model that succeeds only by interpolating among known structures is unlikely to be sufficient, because entire classes of RNA folds are weakly represented or absent from available structural databases. The RNA3DB dataset, for example, was developed to address the sparse and structurally biased nature of RNA structures in the PDB and reports fewer than 2,000 non-redundant examples of RNA structures--two orders of magnitude fewer than are available for proteins \citep{Szikszai2024RNA3DB}.

RNA will therefore be used repeatedly in this thesis as a biomolecular system for evaluating generative models. The goal is not only to generate plausible coordinates, but to learn distributions over structures and how they respond to the environment from sparse data. This motivates the central theme of the thesis: generative models for molecular science require physical and chemical inductive biases, especially when the target system is flexible, data-poor, and ensemble in nature.

\section{The role of physical structure in generative modeling for molecular science}
\label{ch1:sec:physical-structure}

Inductive biases are structural assumptions built into a model that reduce its solution space and, by extension, the amount of data required for learning \citep{Battaglia2018Relational}. In molecular machine learning, these biases take several forms. Some are architectural: translational, rotational, and Euclidean equivariances can be built directly into the network \citep{Fuchs2020SE3}. Others enter through the objective function, such as the addition of physics-informed loss terms that penalize violations of known constraints \citep{Cuomo2022PINNs}. Probabilistic generative models offer a third route: physical structure can be placed directly into the prior distribution or into the dynamical process that transports between the prior and target distribution.

Boltzmann generators are an early example of this strategy \citep{Noe2019Boltzmann}. They learn a map between a simple reference distribution and a molecular Boltzmann distribution while retaining enough structure for reweighting and free-energy estimation. In one reported application, the authors scaled the variance of the prior distribution depending on the temperature at which the training sample was obtained. This enabled the model to generate data at multiple temperatures and reweight samples across temperatures. However, Boltzmann generators require that the energy function be readily available, which limits their application to the full atomistic configuration space. Reduced descriptions instead require a generative model that learns directly from temperature-dependent data, without access to the energy function.

Extrapolation to new environmental conditions cannot be expected to occur solely on the basis of structural data. AlphaFold2 and subsequent AI structure-prediction models are trained largely on structures obtained by X-ray crystallography, NMR spectroscopy, and cryo-electron microscopy \citep{alphafold-2,alphafold-3,berman2000protein}. When such experiments succeed, their output is a molecular structure or structural ensemble corresponding to a particular set of environmental conditions. Structure, however, is linked to the environment: ion concentrations, protein or nucleic-acid partners, cofactors, metabolites, temperature, pressure, and pH can all alter the relevant ensemble. A static structural database offers only glimpses into the full relationship between molecular identity, environment, and ensemble \citep{yamagami2021functional,huang2025some,tiwary2024generative}.

This limitation cannot be overcome simply by collecting the same type of data at every condition of interest. It took decades to assemble the current Protein Data Bank, and it is still sparse for RNAs and heterogeneous biomolecular assemblies \citep{berman2000protein,Szikszai2024RNA3DB}. Repeating high-resolution experiments across every relevant temperature, pressure, ionic condition, ligand concentration, and binding-partner composition would be impractical. Available structures reflect only the set of conditions and limited thermodynamic contexts in which the structure could be resolved.

Statistical mechanics provides a way to approach this problem. The equilibrium distribution describes how temperature changes the weights of configurations in an ensemble. Responses to other thermodynamic variables can be described by introducing the corresponding conjugate variables into the thermodynamic potential. For example, pressure couples to volume, and chemical potential couples to species number or concentration. In either case, the environmental variable changes the probability of conformations through a structured modification to the energy, not an arbitrary transformation.

Rather than expecting such relationships to emerge from data alone, a generative model can be biased to learn them by modifying its probabilistic structure. In a diffusion model, this can be done by changing the prior distribution or the generative dynamics so that the environmental variable is represented explicitly. This modification forms the conceptual basis for Thermodynamic Maps, a core methodology developed throughout the thesis.

\section{Thesis contributions and organization}
\label{ch1:sec:contributions-organization}

This thesis makes three connected contributions. First, it evaluates probabilistic generative models in molecular settings and identifies when particular model classes are useful for distributions with different dimensionality, multimodality, and data scarcity. Second, it develops the idea that thermodynamic structure can be encoded directly into generative models, culminating in the Thermodynamic Maps framework for extrapolating molecular ensembles across temperature and related environmental variables. Third, it translates these principles into practical hybrid workflows: RNAnneal for \textit{ab initio} RNA ensemble prediction from sequence and DiffDock-Glide for molecular docking through a combination of learned proposal generation and physics-based refinement. Across these studies, the common claim is that generative models become more reliable in low-data molecular regimes when their priors, objectives, or pipelines reflect physical knowledge and intuition.

The remainder of the thesis is organized as follows. Chapter~2 benchmarks several probabilistic generative frameworks on molecularly motivated distributions, including Gaussian mixtures with tunable complexity and Aib9 peptide torsion-angle data. These comparisons motivate the later focus on diffusion models for moderate-dimensional, multimodal molecular landscapes. Chapter~3 applies diffusion models to replica-exchange molecular dynamics data and shows that information can be mixed across temperatures to improve conformational sampling and estimate thermodynamic quantities at temperatures not directly simulated.

Chapter~4 introduces Thermodynamic Maps, the core methodological contribution of the thesis. By giving the prior distribution a thermodynamic interpretation, Thermodynamic Maps allow temperature-dependent ensembles and free-energy behavior to be inferred from limited observations. The method is validated on the two-dimensional Ising model and then applied to RNA conformational landscapes through Thermodynamic Map-accelerated molecular dynamics. Chapter~5 builds on this framework in RNAnneal, an \textit{ab initio} pipeline for RNA structure-ensemble prediction that combines hierarchical structure generation, molecular simulation, latent representation learning, and thermodynamic scoring. Chapter~6 applies the broader hybrid physics-AI philosophy to molecular docking in DiffDock-Glide, where a learned diffusion model proposes candidate poses and a physics-based docking engine refines and ranks them. Finally, Chapter~7 discusses extensions to additional thermodynamic variables, latent generative representations, and therapeutically relevant RNA targets.
